% Generated by `taskgen course examples` from number_theory.sqrt_cipolla (seed 5).
% Do not edit: change examples.yaml and regenerate.
\begin{example}[метод Чиполлы]\label{ex:sqrt-cipolla}
Решить квадратичное сравнение методом Чиполлы: \[x^2 \equiv 10\ (\textrm{mod}\ 89).\]

\par\noindent\textit{Решение.} Критерий Эйлера: $10^{(p-1)/2} = 10^{44} \equiv 1$ --- квадратичный вычет. Ищем $t$, для которого $t^2 - 10$ --- невычет (значение $(t^2 - a)^{44}$ равно $-1 \equiv 88$): $t = 1$: $t^2 - a \equiv 80$, $80^{44} \equiv 1$; $t = 2$: $t^2 - a \equiv 83$, $83^{44} \equiv 88$. Положим $\omega^2 = 83$; в $\mathbb{F}_{89}(\omega)$ $(u_1 + v_1\omega)(u_2 + v_2\omega) = (u_1u_2 + 83 v_1v_2) + (u_1v_2 + u_2v_1)\omega$. Вычисляем $x = (2 + \omega)^{(p+1)/2} = (2 + \omega)^{45}$, $45 = 101101_2$, бинарным алгоритмом:\[\begin{array}{c|c|c}\text{бит} & r^2 & r \\ \hline 1 & 1 + 0\omega & 2 + 1\omega \\ 0 & 87 + 4\omega & 87 + 4\omega \\ 1 & 86 + 73\omega & 1 + 54\omega \\ 1 & 38 + 19\omega & 51 + 76\omega \\ 0 & 74 + 9\omega & 74 + 9\omega \\ 1 & 6 + 86\omega & 30 + 0\omega \\ \end{array}\]Получили $x = 30$ (коэффициент при $\omega$ равен $0$), второй корень $p - x = 59$; проверка: $30^2 = 900 \equiv 10$.\[x \equiv \pm 30\ (\textrm{mod}\ 89),\ \text{т.е.}\ x \in \{30, 59\}.\]
\end{example}
